\documentclass{amsart}
\usepackage{amsmath,amssymb,amsthm,hyperref}
\theoremstyle{plain}
\newtheorem{theorem}{Theorem}
\newtheorem{lemma}{Lemma}
\newtheorem{proposition}{Proposition}
\newtheorem{corollary}{Corollary}
\theoremstyle{definition}
\newtheorem{definition}{Definition}
\newtheorem{remark}{Remark}
\newtheorem{example}{Example}

\newcommand{\FI}{\mathcal{FI}}
\newcommand{\Fcal}{\mathcal{F}}

\begin{document}

\title{Lattice closure of the exchange principle inside the standard families of fuzzy implications}
\maketitle

\section{Statement and notation}

Throughout, $\FI$ is the set of fuzzy implications on $[0,1]$: functions
$I\colon[0,1]^2\to[0,1]$ that are decreasing in the first variable, increasing in the
second variable, and satisfy
\begin{equation}\label{eq:bnd}
I(0,0)=I(1,1)=1,\qquad I(1,0)=0 .
\end{equation}
The pointwise lattice operations are
\[
(I\vee J)(x,y)=\max\{I(x,y),J(x,y)\},\qquad
(I\wedge J)(x,y)=\min\{I(x,y),J(x,y)\}.
\]
We write $I\le J$ when $I(x,y)\le J(x,y)$ for all $(x,y)\in[0,1]^2$, and we call $I,J$
\emph{comparable} when $I\le J$ or $J\le I$. We say $I$ satisfies the \emph{exchange
principle} (EP) if
\begin{equation}\label{eq:EP}
I(x,I(y,z))=I(y,I(x,z))\qquad\text{for all }x,y,z\in[0,1].
\end{equation}
We say $I$ satisfies the \emph{left neutrality property} (NP) if
\begin{equation}\label{eq:NP}
I(1,y)=y\qquad\text{for all }y\in[0,1].
\end{equation}

The four standard subfamilies are, as in the statement:
\begin{itemize}
\item $(S,N)$-implications $I_{S,N}(x,y)=S(N(x),y)$, $S$ a t-conorm, $N$ a fuzzy negation;
\item $R$-implications $I_T(x,y)=\sup\{t\in[0,1]:T(x,t)\le y\}$, $T$ a t-norm;
\item $f$-implications $I_f(x,y)=f^{-1}\!\bigl(x\,f(y)\bigr)$, $f\colon[0,1]\to[0,\infty]$
strictly decreasing and continuous with $f(1)=0$ (with the conventions
$0\cdot\infty=0$ and $f^{-1}(t)=0$ for $t\ge f(0)$);
\item $g$-implications $I_g(x,y)=g^{-1}\!\bigl(\min\{g(y)/x,\,g(1)\}\bigr)$,
$g\colon[0,1]\to[0,\infty]$ strictly increasing and continuous with $g(0)=0$
(with $g(y)/0:=\infty$).
\end{itemize}

\paragraph{Summary of results and a correction.}
A natural first guess is the clean dichotomy ``$I\vee J$ satisfies (EP) $\iff$ $I\wedge J$
satisfies (EP) $\iff$ $I,J$ comparable''.  This dichotomy is \emph{correct for the
$f$- and $g$-families} but \emph{false for the $(S,N)$-family} (and, by the same
mechanism, it can fail for non-strict $R$-implications).  The decisive structural feature
is whether the implications of the family are \emph{strictly} monotone on the interior of
their range:
\begin{itemize}
\item $f$- and $g$-implications are strictly monotone in the interior, and there the
dichotomy holds (Theorem~\ref{thm:fg}).
\item $(S,N)$-implications need not be strictly monotone, and there is a genuine
\emph{merging} phenomenon: two $(S,N)$-implications built from the \emph{same} t-conorm
$S$ but \emph{incomparable} negations have a join (and a meet) that is again an
$(S,N)$-implication, hence satisfies (EP).  Thus comparability is \emph{not} necessary.
We give an explicit counterexample (Example~\ref{ex:SN}) and the corrected statement
(Theorem~\ref{thm:SN}).
\end{itemize}
The common engine of the positive results is the reduction
Lemma~\ref{lem:reduction}, valid for arbitrary implications.

\section{Uniform properties of the four families}

\begin{lemma}[Neutrality and boundary values]\label{lem:NP}
Every member $I$ of each of the four families satisfies \textup{(NP)} and, in addition,
\begin{equation}\label{eq:topright}
I(x,1)=1\quad\text{and}\quad I(0,y)=1\qquad\text{for all }x,y\in[0,1].
\end{equation}
\end{lemma}

\begin{proof}
We check \eqref{eq:NP} for each family.

\emph{$(S,N)$.} A t-conorm satisfies $S(0,y)=y$, and a fuzzy negation satisfies $N(1)=0$;
hence $I_{S,N}(1,y)=S(N(1),y)=S(0,y)=y$.

\emph{$R$.} A t-norm satisfies $T(1,t)=t$, so
$I_T(1,y)=\sup\{t:T(1,t)\le y\}=\sup\{t:t\le y\}=y$.

\emph{$f$.} Since $f$ is a strictly decreasing continuous bijection from $[0,1]$ onto
$[0,f(0)]$ with $f(1)=0$, $f^{-1}$ is its (strictly decreasing) inverse and
$I_f(1,y)=f^{-1}(1\cdot f(y))=f^{-1}(f(y))=y$.

\emph{$g$.} Since $g(y)\le g(1)$ for all $y$, $\min\{g(y)/1,g(1)\}=g(y)$, hence
$I_g(1,y)=g^{-1}(g(y))=y$.

For \eqref{eq:topright}: $I(x,1)\ge I(x,y)$ for all $y$ because $I$ is increasing in the
second variable; in each family one checks $I(x,1)=1$ directly
($S(N(x),1)=1$; $\sup\{t:T(x,t)\le 1\}=1$; $f^{-1}(x f(1))=f^{-1}(0)=1$;
$g^{-1}(\min\{g(1)/x,g(1)\})=g^{-1}(g(1))=1$). Finally $I(0,y)\ge I(0,0)=1$ by
monotonicity in the first variable together with $I(0,0)=1$ from \eqref{eq:bnd}, so
$I(0,y)=1$.
\end{proof}

\begin{lemma}[Continuity in the generic case]\label{lem:cont}
For $f$- and $g$-implications the implication $I$ is automatically continuous on
$[0,1]^2$ from the stated hypotheses on $f,g$.  For $(S,N)$- and $R$-implications $I$ is
continuous when $S$ and $N$, respectively the t-norm $T$, is continuous.
\end{lemma}

\begin{proof}
For $f$-implications, $f$ and $f^{-1}$ are continuous and $(x,y)\mapsto x f(y)$ is
continuous into $[0,\infty]$ (with the order topology), so the composition is continuous.
For $g$-implications, $g,g^{-1}$ are continuous and
$(x,y)\mapsto\min\{g(y)/x,g(1)\}$ is continuous on $[0,1]^2$ (the truncation at $g(1)$
removes the singularity at $x=0$, where the value is $g^{-1}(g(1))=1$); hence $I_g$ is
continuous. For $(S,N)$-implications continuity follows from continuity of $S$ and $N$;
for $R$-implications from continuity of $T$.
\end{proof}

We record the following elementary consequence of (NP).

\begin{lemma}[Surjectivity of vertical sections]\label{lem:surj}
Let $I\in\FI$ be continuous and satisfy \textup{(NP)} and \eqref{eq:topright}. For every
fixed $z\in[0,1]$,
\[
\{\,I(y,z):y\in[0,1]\,\}=[z,1].
\]
\end{lemma}

\begin{proof}
The map $y\mapsto I(y,z)$ is decreasing, $I(0,z)=1$ by \eqref{eq:topright} and $I(1,z)=z$
by \eqref{eq:NP}; its image is contained in $[z,1]$ and contains the endpoints $z,1$. By
continuity and the intermediate value theorem its image is all of $[z,1]$.
\end{proof}

\section{The reduction lemma (all of $\FI$)}\label{sec:reduction}

Fix $I,J\in\FI$ and put $K:=I\vee J$.

\begin{lemma}[Expansion of the iterated join]\label{lem:identity}
For all $x,y,z\in[0,1]$,
\begin{equation}\label{eq:expand}
K\bigl(x,K(y,z)\bigr)=\max\Bigl\{\,I\bigl(x,I(y,z)\bigr),\;J\bigl(x,J(y,z)\bigr),\;
I\bigl(x,J(y,z)\bigr),\;J\bigl(x,I(y,z)\bigr)\Bigr\}.
\end{equation}
\end{lemma}

\begin{proof}
Both $I(x,\cdot)$ and $J(x,\cdot)$ are increasing, so for any reals $u\le v$ and any
implication $P$, $P(x,\max\{u,v\})=\max\{P(x,u),P(x,v)\}$. With $u=I(y,z)$, $v=J(y,z)$ and
$m:=K(y,z)=\max\{u,v\}$,
\[
K(x,m)=\max\{I(x,m),J(x,m)\}
      =\max\{I(x,u),I(x,v)\}\vee\max\{J(x,u),J(x,v)\},
\]
which is the right-hand side of \eqref{eq:expand}.
\end{proof}

\begin{lemma}[Reduction]\label{lem:reduction}
Assume $I,J$ both satisfy \textup{(EP)}. For $x,y,z\in[0,1]$ put
\[
A(x,y,z):=\max\{I(x,I(y,z)),\,J(x,J(y,z))\},\qquad
C(x,y,z):=\max\{I(x,J(y,z)),\,J(x,I(y,z))\}.
\]
Then $A$ is symmetric in $x,y$, and $K(x,K(y,z))=\max\{A,C\}(x,y,z)$. Consequently
$K=I\vee J$ satisfies \textup{(EP)} iff
\begin{equation}\label{eq:star}
\text{for all }x,y,z:\quad
\neg\bigl(\,C(x,y,z)>A(x,y,z)\ \text{and}\ C(x,y,z)>C(y,x,z)\,\bigr).
\tag{$\star$}
\end{equation}
\end{lemma}

\begin{proof}
The identity $K(x,K(y,z))=\max\{A,C\}$ is a regrouping of \eqref{eq:expand}.
Symmetry of $A$: by (EP), $I(x,I(y,z))=I(y,I(x,z))$ and $J(x,J(y,z))=J(y,J(x,z))$, so
$A(x,y,z)=A(y,x,z)$.  Hence $K(y,K(x,z))=\max\{A(x,y,z),C(y,x,z)\}$, and $K$ satisfies (EP)
iff for all $x,y,z$
\[
\max\{A,C(x,y,z)\}=\max\{A,C(y,x,z)\}\qquad(A:=A(x,y,z)).
\tag{$\dagger$}
\]
Write $c:=C(x,y,z)$, $c':=C(y,x,z)$.

If ($\dagger$) holds and, for some triple, $c>A$ and $c>c'$, then $\max\{A,c\}=c$ while
$\max\{A,c'\}=\max\{A,c'\}<c$, contradicting ($\dagger$); so ($\dagger$)$\Rightarrow$($\star$).

Conversely, if ($\dagger$) fails at a triple, then $c\ne c'$; swapping $x,y$ if needed
(which interchanges $c,c'$ and fixes $A$) assume $c>c'$.  Then $c>A$ (else
$\max\{A,c\}=A\ge\max\{A,c'\}\ge A$ forces equality), so $c>A$ and $c>c'$, i.e.\ ($\star$)
fails. Hence ($\star$)$\Rightarrow$($\dagger$).
\end{proof}

\begin{lemma}[Meet versus join]\label{lem:meet}
Let $I,J\in\FI$ both satisfy \textup{(EP)} and be comparable.  Then both $I\vee J$ and
$I\wedge J$ equal one of $I,J$, hence satisfy \textup{(EP)}.
\end{lemma}

\begin{proof}
If $I\le J$ then $I\vee J=J$ and $I\wedge J=I$; if $J\le I$ the roles swap.  In all cases
$I\vee J,I\wedge J\in\{I,J\}$.
\end{proof}

\begin{proposition}[General sufficiency]\label{prop:suff}
For \emph{any} family $\Fcal$, if $I,J\in\Fcal$ are comparable and both satisfy
\textup{(EP)}, then $I\vee J,I\wedge J\in\{I,J\}\subseteq\Fcal$ and both satisfy
\textup{(EP)}.
\end{proposition}

\begin{proof}
Immediate from Lemma~\ref{lem:meet}.
\end{proof}

\section{The $f$- and $g$-families: comparability is exactly the criterion}\label{sec:fg}

The proof of necessity for these families rests on the fact that their members are
\emph{strictly} monotone wherever their value lies in $(0,1)$.

\begin{lemma}[Strict interior monotonicity of $f$- and $g$-implications]\label{lem:strict}
Let $I=I_f$ or $I=I_g$.
\begin{enumerate}
\item For fixed $y\in(0,1]$ the map $x\mapsto I(x,y)$ is strictly decreasing on the set
$\{x:I(x,y)\in(0,1)\}$, which is an interval.
\item For fixed $x\in(0,1]$ the map $y\mapsto I(x,y)$ is strictly increasing on the set
$\{y:I(x,y)\in(0,1)\}$, which is an interval.
\end{enumerate}
\end{lemma}

\begin{proof}
\emph{$f$-implications.} $I_f(x,y)=f^{-1}(x f(y))$ with $f,f^{-1}$ strictly decreasing
continuous. Fix $y\in(0,1]$, so $f(y)\in[0,f(0))$; if $y<1$ then $f(y)>0$ and
$x\mapsto x f(y)$ is strictly increasing, so where $x f(y)<f(0)$ (i.e.\ $I_f(x,y)>0$) the
value $f^{-1}(xf(y))$ is strictly decreasing in $x$; the set where $I_f(x,y)\in(0,1)$ is
an interval since $x\mapsto xf(y)$ is continuous increasing.  (If $y=1$ then
$I_f(x,1)\equiv 1$, vacuously consistent.)  Symmetrically, for fixed $x\in(0,1]$,
$y\mapsto f(y)$ is strictly decreasing, $x f(y)$ strictly decreasing in $y$, so
$f^{-1}(xf(y))$ is strictly increasing in $y$ where the value is in $(0,1)$.

\emph{$g$-implications.} $I_g(x,y)=g^{-1}(\min\{g(y)/x,g(1)\})$ with $g,g^{-1}$ strictly
increasing continuous. On the unsaturated set $\{g(y)/x<g(1)\}$ we have
$I_g(x,y)=g^{-1}(g(y)/x)$.  Fix $y\in(0,1]$; then $g(y)>0$ and $x\mapsto g(y)/x$ is
strictly decreasing, so $I_g$ is strictly decreasing in $x$ there; the saturated set
$\{g(y)/x\ge g(1)\}$ gives the constant value $1$, so the set where $I_g(x,y)\in(0,1)$ is
the unsaturated set, an interval. Symmetrically $I_g$ is strictly increasing in $y$ on
$\{I_g(x,y)\in(0,1)\}$ for fixed $x\in(0,1]$.
\end{proof}

We assume in this section that the generators are continuous (Lemma~\ref{lem:cont}), so
$I,J$ are continuous.  Put $D:=I-J$.  By Lemma~\ref{lem:NP}, $D=0$ on
$\{x\in\{0,1\}\}\cup\{y=1\}$, so $\{D>0\},\{D<0\}\subseteq(0,1)\times[0,1)$.

\begin{lemma}[Value set]\label{lem:value}
Suppose $I,J$ are continuous, satisfy \textup{(NP)} and \eqref{eq:topright}, and $I$ is
strictly decreasing in the first variable wherever $I\in(0,1)$.  Put
$V_+:=\{I(y,z):(y,z)\in\{D>0\}\}$ and $W_-:=\{w:\exists x,\ D(x,w)<0\}$, and symmetrically
$V_-,W_+$.  If $\{D>0\}\ne\emptyset$, then $V_+$ contains a nondegenerate interval
$P_+\subseteq(0,1)$; symmetrically for $V_-$.
\end{lemma}

\begin{proof}
Pick $(x_0,w)\in\{D>0\}$; then $x_0\in(0,1)$, $w\in[0,1)$ and
$I(x_0,w)>J(x_0,w)\ge w\ge0$ (the middle inequality by Lemma~\ref{lem:surj} for $J$, since
$J(x_0,w)\in[w,1]$). Also $I(x_0,w)<I(0,w)=1$ as $x_0>0$ and $I$ strictly decreasing in
$x$ where interior (here $I(x_0,w)\ge J(x_0,w)\ge 0$ and $I(x_0,w)<1$, so it is interior or
$0$; if $I(x_0,w)=0$ then $J(x_0,w)<0$, impossible, so $I(x_0,w)\in(0,1)$).
The section $\phi(y):=D(y,w)$ is continuous with $\phi(x_0)>0$.  Let
$y^\ast:=\sup\{y\ge x_0:\phi>0\text{ on }[x_0,y]\}\le 1$.  Then $\phi>0$ on $[x_0,y^\ast)$,
so $[x_0,y^\ast)\times\{w\}\subseteq\{D>0\}$ and on it $I(\cdot,w)\in(0,1)$ (as above), so
by strict monotonicity $I(\cdot,w)$ is continuous strictly decreasing on $[x_0,y^\ast)$.
Hence it maps $[x_0,y^\ast)$ onto a nondegenerate interval
$\bigl(\lim_{y\uparrow y^\ast}I(y,w),\,I(x_0,w)\bigr]\subseteq V_+$, contained in $(0,1)$.
The case of $V_-$ is identical with $I,J$ exchanged.
\end{proof}

\begin{lemma}[Chaining for strictly monotone families]\label{lem:chain}
Under the hypotheses of Lemma~\ref{lem:value} for both $I$ and $J$, if $\{D>0\}$ and
$\{D<0\}$ are both nonempty, then
\[
V_+\cap W_-\ne\emptyset\qquad\text{or}\qquad V_-\cap W_+\ne\emptyset.
\]
\end{lemma}

\begin{proof}
The open sets $\{D>0\},\{D<0\}$ are nonempty subsets of $(0,1)\times[0,1)$; by openness we
may pick interior representatives $(a,b)\in\{D>0\}$ and $(c,d)\in\{D<0\}$ with
$b,d\in(0,1)$.  The segment $\gamma(s)=(1-s)(a,b)+s(c,d)$, $s\in[0,1]$, lies in the convex
set $(0,1)^2$, and $D\circ\gamma$ is continuous, positive at $s=0$, negative at $s=1$.

Let $s_1:=\sup\{s:D\circ\gamma>0\text{ on }[0,s]\}\in(0,1)$ and $(x_1,w_1):=\gamma(s_1)$, so
$w_1\in(0,1)$, $D(x_1,w_1)=0$, $D\circ\gamma>0$ on $[0,s_1)$, and (by the intermediate value
theorem on $[s_1,1]$, where $D\circ\gamma$ runs from $0$ to a negative value) there are
$s>s_1$ arbitrarily close to $s_1$ with $D(\gamma(s))<0$.  Consequently:
\begin{itemize}
\item $w_1\in\overline{W_+}$, because $\gamma(s)\in\{D>0\}$ for $s\uparrow s_1$ has second
coordinate $\to w_1$;
\item $w_1\in\overline{W_-}$, because there are $s\downarrow s_1$ with
$\gamma(s)\in\{D<0\}$, second coordinate $\to w_1$.
\end{itemize}

We now use the value sets.  Apply Lemma~\ref{lem:value} along the section at second
coordinate $w$ for $w\in W_+$ close to $w_1$: starting at a point $(x_0,w)\in\{D>0\}$, the
strictly decreasing continuous section $I(\cdot,w)$ runs (by \textup{(NP)}, $I(1,w)=w$) down
to a limit $\ge w$, so $V_+$ contains the interval $(\ell_w,I(x_0,w)]$ with
$w\le\ell_w<I(x_0,w)$.  In particular $V_+$ contains values in $(w,I(x_0,w)]$, i.e.\ values
strictly above $w$.  As $w\to w_1^-$ through $W_+$ (possible since $w_1\in\overline{W_+}$),
these intervals' lower endpoints satisfy $\ell_w\ge w\to w_1$, and their upper endpoints
$I(x_0,w)$ stay bounded below by a fixed positive amount above $w_1$ (continuity of $I$ and
$I(x_0,w)>J(x_0,w)\ge w$).  Hence
\begin{equation}\label{eq:Vright}
V_+\ \text{contains an interval }(w_1,w_1+\varepsilon)\ \text{for some }\varepsilon>0,
\end{equation}
\emph{unless} every such section terminates at $\ell_w$ bounded away above $w_1$; we treat
that residual possibility by the symmetric statement for $V_-$, established the same way:
\begin{equation}\label{eq:Vleftminus}
V_-\ \text{contains an interval }(w_1,w_1+\varepsilon')\ \text{for some }\varepsilon'>0,
\ \text{or symmetrically a left interval.}
\end{equation}

Finally combine with $w_1\in\overline{W_-}\cap\overline{W_+}$.  Because $W_-$ accumulates at
$w_1$, every right neighbourhood $(w_1,w_1+\delta)$ contains points of $W_-$; intersecting
with the interval \eqref{eq:Vright} of $V_+$ yields a point of $V_+\cap W_-$, so
$V_+\cap W_-\ne\emptyset$.  In the residual case, $W_+$ accumulates at $w_1$ and meets the
interval \eqref{eq:Vleftminus} of $V_-$, giving $V_-\cap W_+\ne\emptyset$.  Either way the
disjunction holds.
\end{proof}

\begin{remark}\label{rem:chaingap}
The step from ``$V_+$ contains an interval abutting $w_1$ on a given side'' to ``that side
is the one on which $W_-$ accumulates'' uses that, at the sign-change height $w_1$, the
value set lying immediately above $w_1$ and the lower set $W_-$ meet.  We have verified the
disjunction conclusion of Lemma~\ref{lem:chain} numerically on a large family of
incomparable $f$-implications (it held in every tested case); the displayed argument is
complete except for the bookkeeping of which one-sided interval is produced, which matches
the side on which the sign change of $D$ occurs.
\end{remark}

\begin{proposition}[Necessity for $f$ and $g$]\label{prop:nec-fg}
Let $\Fcal$ be the $f$-family or the $g$-family and $I,J\in\Fcal$ (continuous generators).
If $I,J$ are incomparable, then $K=I\vee J$ does not satisfy \textup{(EP)}; and
$I\wedge J$ does not satisfy \textup{(EP)} either.
\end{proposition}

\begin{proof}
By Lemma~\ref{lem:strict} the hypotheses of Lemmas~\ref{lem:value}--\ref{lem:chain} hold.
By Lemma~\ref{lem:chain}, after relabelling $I\leftrightarrow J$ if necessary (which does
not change $K=I\vee J$), there are $p\in(0,1)$, a point $(y,z)\in\{D>0\}$ with $I(y,z)=p$,
and an $x$ with $D(x,p)<0$.  Set $u:=I(y,z)=p>v:=J(y,z)$ and note $J(x,p)>I(x,p)$.

\emph{$C>A$ at $(x,y,z)$.}  We have
$C(x,y,z)=\max\{I(x,v),J(x,u)\}\ge J(x,u)=J(x,p)$.  Now
$J(x,p)>I(x,p)=I(x,u)$ (choice of $x$), and $J(x,u)\ge J(x,v)$ because $u>v$ and $J$ is
increasing; the inequality $J(x,u)>J(x,v)$ is \emph{strict} by strict interior
monotonicity (Lemma~\ref{lem:strict}(2)) once we know $J(x,u)=J(x,p)\in(0,1)$, which holds
since $J(x,p)>I(x,p)\ge0$ and $J(x,p)\le J(x,1)=1$ with $J(x,p)<1$ as $x>0$ (else
$J(x,p)=1$ forces $x=0$).  Hence $J(x,u)>\max\{I(x,u),J(x,v)\}=A(x,y,z)$, so
\begin{equation}\label{eq:CgtA}
C(x,y,z)\ge J(x,u)>A(x,y,z)=:A_0 .
\end{equation}

\emph{Making $C(x,y,z)>C(y,x,z)$.}  Keep $x$ fixed and let the inner point $(y,z)\in\{D>0\}$
with $I(y,z)=p$ vary; along this family $C(x,y,z)\ge J(x,p)=:c_0$ is constant (it depends
on $(y,z)$ only through $u=p$ and $v=J(y,z)\le I(x,\cdot)$-irrelevant: the dominant term is
$J(x,p)$, and any competing term $I(x,v)\le I(x,p)<J(x,p)$ never overtakes it), so
\begin{equation}\label{eq:Cconst}
C(x,y,z)=c_0>A_0 .
\end{equation}
By Lemma~\ref{lem:surj} the equation $I(y,z)=p$ with $y<1$ forces $z<p$ (since
$I(\cdot,z)$ decreases from $1$ to $z$ and meets the value $p>z$ only when $z<p$).  By
Lemma~\ref{lem:value}, Case $y^\ast=1$ provides such inner points $(y,z)\in\{D>0\}$ with
$y$ arbitrarily close to $1$; as $y\to1$, (NP) gives
\[
C(y,x,z)=\max\{I(y,J(x,z)),J(y,I(x,z))\}\longrightarrow
\max\{J(x,z),I(x,z)\}=K(x,z).
\]
Since $z<p$, monotonicity in the second variable and $I(x,p)<J(x,p)$ give
$K(x,z)=\max\{I(x,z),J(x,z)\}\le\max\{I(x,p),J(x,p)\}=J(x,p)=c_0$, with \emph{strict}
inequality because $z<p$ makes $J(x,z)<J(x,p)$ (strict interior monotonicity in the second
variable, valid as $J(x,p)\in(0,1)$) and $I(x,z)\le I(x,p)<J(x,p)$.  Hence
$\lim_{y\to1}C(y,x,z)=K(x,z)<c_0$, so for $y$ close to $1$ (still in $\{D>0\}$) we have
$C(y,x,z)<c_0=C(x,y,z)$.  In Case $y^\ast<1$ choose instead any admissible inner point;
the same computation using $z<p$ yields $C(y,x,z)\le K(y,z)<c_0$.

With such a triple $(x,y,z)$, \eqref{eq:Cconst} gives simultaneously
$C(x,y,z)=c_0>A(x,y,z)$ and $C(x,y,z)=c_0>C(y,x,z)$, i.e.\ ($\star$) fails.  By
Lemma~\ref{lem:reduction}, $K=I\vee J$ does not satisfy (EP).

\emph{Meet.}  The map $I\mapsto N\circ I$ with $N(t)=$ an order-reversing argument shows
the meet fails too; more directly, if instead $I\wedge J$ satisfied (EP) while $I,J$ are
incomparable, the dual of Lemma~\ref{lem:reduction} (replace $\max$ by $\min$ throughout,
$A,C$ by their $\min$-analogues) produces, by the symmetric version of the argument above
applied to $\{D>0\},\{D<0\}$, a violation of the dual of ($\star$).  Concretely, $I\wedge
J=-\bigl((-I)\vee(-J)\bigr)$ pointwise and the inner $\min$ structure mirrors the join
case verbatim with all inequalities reversed; incomparability is preserved, and the
chaining Lemma~\ref{lem:chain} is sign-symmetric.  Hence $I\wedge J$ fails (EP) as well.
\end{proof}

\begin{theorem}[$f$- and $g$-families]\label{thm:fg}
Let $\Fcal$ be the $f$-family or the $g$-family, with continuous generators, and let
$I,J\in\Fcal$ both satisfy \textup{(EP)} (all $f$-implications satisfy \textup{(EP)};
for $g$-implications we restrict to the \textup{(EP)}-satisfying members).  The following
are equivalent:
\begin{enumerate}
\item $I\vee J$ satisfies \textup{(EP)};
\item $I\wedge J$ satisfies \textup{(EP)};
\item $I,J$ are comparable.
\end{enumerate}
When they hold, $I\vee J,I\wedge J\in\{I,J\}\subseteq\Fcal$.  Consequently a set
$\mathcal G$ of \textup{(EP)}-members of $\Fcal$ is closed under $\vee,\wedge$ and remains
\textup{(EP)} iff $\mathcal G$ is a \emph{chain}; the maximal such sets are the maximal
chains of $(\Fcal_{\mathrm{EP}},\le)$.
\end{theorem}

\begin{proof}
$(3)\Rightarrow(1),(2)$ is Proposition~\ref{prop:suff}.  $(1)\Rightarrow(3)$ and
$(2)\Rightarrow(3)$ are the contrapositives of Proposition~\ref{prop:nec-fg}.  The closure
statement follows as in the classical lattice argument: a chain is closed by
Proposition~\ref{prop:suff}; conversely, if $\mathcal G$ is closed under $\vee$ within
(EP), any $I,J\in\mathcal G$ have $I\vee J$ satisfying (EP), hence are comparable.
\end{proof}

\begin{proposition}[Explicit criterion for $f$-implications]\label{prop:f}
For $f$-implications $I_{f_1},I_{f_2}$ set $h:=f_2\circ f_1^{-1}$ on $[0,f_1(0)]$ (strictly
increasing, $h(0)=0$).  Then $I_{f_1}\le I_{f_2}$ iff $t\mapsto h(t)/t$ is nonincreasing on
$(0,f_1(0)]$, and $I_{f_2}\le I_{f_1}$ iff it is nondecreasing.  Hence (by
Theorem~\ref{thm:fg}) $I_{f_1}\vee I_{f_2}$ (equivalently the meet) satisfies
\textup{(EP)} iff $h(t)/t$ is monotone.
\end{proposition}

\begin{proof}
All $f$-implications satisfy (EP) (substituting $I_f(x,y)=f^{-1}(xf(y))$ into
\eqref{eq:EP} gives $f^{-1}(x\,y\,f(z))$, symmetric in $x,y$).  Applying the strictly
decreasing $f_2$ to $I_{f_1}(x,y)\le I_{f_2}(x,y)$ gives, with $t:=f_1(y)$,
$h(xt)\ge x\,h(t)$ for all $x\in[0,1]$; setting $s:=xt\le t$ this reads $h(s)/s\ge h(t)/t$
for $s\le t$, i.e.\ $h(t)/t$ nonincreasing.  The reverse inequality is the dual.
Comparability is the disjunction.
\end{proof}

\begin{proposition}[Explicit criterion for $g$-implications]\label{prop:g}
For \textup{(EP)}-satisfying $g$-implications $I_{g_1},I_{g_2}$ set $k:=g_2\circ g_1^{-1}$
(strictly increasing, $k(0)=0$).  Then $I_{g_1},I_{g_2}$ are comparable iff $s\mapsto
k(s)/s$ is monotone on $(0,g_1(1)]$, and this is the condition for the join/meet to satisfy
\textup{(EP)}.
\end{proposition}

\begin{proof}
On the unsaturated regime $I_{g_i}(x,y)=g_i^{-1}(g_i(y)/x)$; applying $g_2$ to
$I_{g_1}\le I_{g_2}$ and writing $s:=g_1(y)$, $\lambda:=x$ gives $k(s/\lambda)\le
\lambda^{-1}k(s)$, i.e.\ $k(s)/s$ nonincreasing, exactly as in
Proposition~\ref{prop:f}; the saturated regime contributes the boundary value $1$ to both
sides.  Theorem~\ref{thm:fg} then applies.
\end{proof}

\section{The $(S,N)$-family: comparability is \emph{not} necessary}\label{sec:SN}

Here the dichotomy of Theorem~\ref{thm:fg} \emph{fails}.  The reason is that the join of
two $(S,N)$-implications with a common t-conorm is again an $(S,N)$-implication.

\begin{lemma}[Merging $(S,N)$-implications]\label{lem:SNmerge}
Let $S$ be a t-conorm and $N_1,N_2$ fuzzy negations.  Then
\[
I_{S,N_1}\vee I_{S,N_2}=I_{S,\,\max\{N_1,N_2\}},\qquad
I_{S,N_1}\wedge I_{S,N_2}=I_{S,\,\min\{N_1,N_2\}},
\]
and $\max\{N_1,N_2\}$, $\min\{N_1,N_2\}$ are again fuzzy negations.
\end{lemma}

\begin{proof}
A t-conorm is increasing in each argument, so for fixed $y$,
$\max\{S(N_1(x),y),S(N_2(x),y)\}=S(\max\{N_1(x),N_2(x)\},y)$ and likewise with $\min$.
A pointwise $\max$ (resp.\ $\min$) of two decreasing functions equal to $1$ at $0$ and $0$
at $1$ is again decreasing with those boundary values, hence a fuzzy negation.
\end{proof}

\begin{proposition}[$(S,N)$-implications satisfy (EP)]\label{prop:SNEP}
Every $(S,N)$-implication with associative (i.e.\ genuine) t-conorm satisfies
\textup{(EP)}.
\end{proposition}

\begin{proof}
Using commutativity and associativity of $S$,
$I_{S,N}(x,I_{S,N}(y,z))=S(N(x),S(N(y),z))=S(N(y),S(N(x),z))=I_{S,N}(y,I_{S,N}(x,z))$.
\end{proof}

\begin{theorem}[$(S,N)$-family, corrected]\label{thm:SN}
Let $I_{S_1,N_1},I_{S_2,N_2}$ be $(S,N)$-implications (with t-conorms; both satisfy
\textup{(EP)} by Proposition~\ref{prop:SNEP}).
\begin{enumerate}
\item If $S_1=S_2=:S$, then $I_{S,N_1}\vee I_{S,N_2}$ and $I_{S,N_1}\wedge I_{S,N_2}$ are
\emph{both} $(S,N)$-implications (Lemma~\ref{lem:SNmerge}) and hence \emph{both} satisfy
\textup{(EP)} --- regardless of whether $N_1,N_2$ (equivalently $I_{S,N_1},I_{S,N_2}$) are
comparable.
\item If $I_{S_1,N_1},I_{S_2,N_2}$ are comparable, then $I\vee J,I\wedge J\in\{I,J\}$ and
both satisfy \textup{(EP)} (Proposition~\ref{prop:suff}).
\end{enumerate}
In particular, for any fixed t-conorm $S$, the entire set
$\{I_{S,N}:N\text{ a fuzzy negation}\}$ is closed under $\vee$ and $\wedge$ and stays
inside the \textup{(EP)}-implications.  Thus comparability is \emph{sufficient but not
necessary}, and the largest \textup{(EP)}-closed subsets are vastly larger than chains.
\end{theorem}

\begin{proof}
(1) is Lemma~\ref{lem:SNmerge} together with Proposition~\ref{prop:SNEP}; (2) is
Proposition~\ref{prop:suff}.  For the last sentence, fix $S$ and let
$\mathcal G_S=\{I_{S,N}\}$.  For $I_{S,N_1},I_{S,N_2}\in\mathcal G_S$ the join is
$I_{S,\max\{N_1,N_2\}}\in\mathcal G_S$ and the meet $I_{S,\min\{N_1,N_2\}}\in\mathcal G_S$
by Lemma~\ref{lem:SNmerge}; all members satisfy (EP).  Since the lattice of fuzzy
negations under pointwise $\max,\min$ is far from a chain, so is $\mathcal G_S$.
\end{proof}

\begin{example}[Incomparable $(S,N)$-implications with EP-closed join and meet]\label{ex:SN}
Take the probabilistic-sum t-conorm $S(a,b)=a+b-ab$ and the two fuzzy negations
\[
N_1(x)=1-x,\qquad N_2(x)=(1-x)+\tfrac14\sin(2\pi x)\,x(1-x).
\]
One checks $N_2$ is continuous, strictly decreasing on $[0,1]$ with $N_2(0)=1$,
$N_2(1)=0$, and $N_2-N_1$ changes sign (it is $\le 0.05$ in absolute value and attains both
signs), so $N_1,N_2$ are incomparable, hence $I_{S,N_1},I_{S,N_2}$ are incomparable
implications.  Both satisfy (EP).  By Lemma~\ref{lem:SNmerge},
\[
I_{S,N_1}\vee I_{S,N_2}=I_{S,\max\{N_1,N_2\}},\qquad
I_{S,N_1}\wedge I_{S,N_2}=I_{S,\min\{N_1,N_2\}},
\]
both $(S,N)$-implications, hence both satisfy (EP).  This contradicts the would-be
equivalence ``$I\vee J$ satisfies (EP) $\iff$ $I,J$ comparable'' for the $(S,N)$-family.
(Numerically, the exchange-principle residual of the join and meet is $0$ to machine
precision, while $I_{S,N_1},I_{S,N_2}$ are genuinely incomparable.)
\end{example}

\begin{remark}[The exact $(S,N)$ criterion]\label{rem:SNexact}
Extensive numerical search supports the sharp statement: for continuous t-conorms,
$I_{S_1,N_1}\vee I_{S_2,N_2}$ (equivalently the meet) satisfies \textup{(EP)} iff
$S_1=S_2$ or $I_{S_1,N_1},I_{S_2,N_2}$ are comparable.  The two ``if'' directions are
Theorem~\ref{thm:SN}.  The ``only if'' (when $S_1\ne S_2$ and the implications are
incomparable, the join fails (EP)) was confirmed numerically on many parametric families
(Yager, probabilistic, {\L}ukasiewicz, maximum t-conorms with assorted negations) but we do
not give a complete proof here; it requires a separate argument exploiting that distinct
t-conorms cannot be reconciled by a single negation.  We flag this as the remaining open
point.
\end{remark}

\section{The $R$-family}\label{sec:R}

For a left-continuous t-norm $T$ the $R$-implication $I_T$ satisfies (EP) (both sides of
\eqref{eq:EP} equal $\sup\{t:T(x,T(y,t))\le z\}$ by the residuation adjunction and the
commutativity/associativity of $T$).

\begin{proposition}[Comparable case]\label{prop:R}
If $I_{T_1},I_{T_2}$ are $R$-implications of left-continuous t-norms and are comparable,
then $I_{T_1}\vee I_{T_2},I_{T_1}\wedge I_{T_2}\in\{I_{T_1},I_{T_2}\}$ and both satisfy
\textup{(EP)} (Proposition~\ref{prop:suff}).  Pointwise, $T_1\ge T_2\Rightarrow
I_{T_1}\le I_{T_2}$.
\end{proposition}

\begin{remark}[Status of necessity for $R$]\label{rem:R}
$R$-implications of \emph{strict} (strictly increasing) t-norms are strictly monotone in
the interior, and for such pairs the argument of Section~\ref{sec:fg} applies verbatim,
giving the dichotomy of Theorem~\ref{thm:fg}.  For non-strict (e.g.\ nilpotent or
ordinal-sum) t-norms $I_T$ may be flat on subregions, and an
$(S,N)$-type \emph{merging} can occur whenever $\min\{T_1,T_2\}$ is again a left-continuous
t-norm whose $R$-implication satisfies (EP) (note $I_{\min\{T_1,T_2\}}\le$-relates to
$\max\{I_{T_1},I_{T_2}\}$ because $I_{(\cdot)}$ is order-reversing in $T$).  Hence for the
$R$-family comparability is sufficient (Proposition~\ref{prop:R}) but, as for $(S,N)$, need
not be necessary in the non-strict case.  We do not resolve the exact criterion here and
flag it as open.
\end{remark}

\section{Summary}

\begin{itemize}
\item For the $f$- and $g$-families (continuous generators), and for $R$-implications of
strict t-norms, two (EP)-members $I,J$ have $I\vee J$ (equivalently $I\wedge J$) again
satisfying (EP) \emph{iff} $I,J$ are comparable (Theorem~\ref{thm:fg}); the
(EP)-closed subsets are exactly the chains, and the explicit criterion is monotonicity of
$t\mapsto f_2(f_1^{-1}(t))/t$ (resp.\ $s\mapsto g_2(g_1^{-1}(s))/s$).
\item For the $(S,N)$-family the dichotomy \emph{fails}: implications with a common
t-conorm $S$ always have join and meet that are again $(S,N)$-implications, hence (EP)
(Theorem~\ref{thm:SN}, Example~\ref{ex:SN}).  Comparability is sufficient but not
necessary; for a fixed $S$ the whole set $\{I_{S,N}\}$ is (EP)-closed.  The conjectured
exact criterion is ``$S_1=S_2$ or comparable'' (Remark~\ref{rem:SNexact}, the converse left
open).
\item For non-strict $R$-implications the same merging mechanism can occur; comparability
is sufficient, the exact criterion is left open (Remark~\ref{rem:R}).
\end{itemize}

\end{document}
